\SourceLeaf{385}

\par\bigskip\Needspace{8\baselineskip}
\pdfbookmark[1]{Bemerkungen zur Theorie der biquadratischen Reste}{nachlass-biquadratic-bemerkungen}
\begin{center}\large\scshape Bemerkungen.\end{center}
\BookRule{.2\textwidth}
Nach der Stelle zu urteilen, die die betreffenden Handschriften in einem Notizbuch einnehmen, stammen die hier im Druck mit I und II bezeichneten Bruchstücke aus dem Jahr $1811$ oder der unmittelbar folgenden Zeit. Von den vorangegangenen Versuchen, das Fundamentaltheorem für biquadratische Reste mit den hier für den Rest $1+i$ angewandten Methoden zu beweisen, ist eine Aufzeichnung erhalten, die den besonderen Fall des Restes $1+2i$ erledigt und von jener Bestimmung des biquadratischen Charakters ausgeht, die Artikel $6$ des Bruchstücks III als Anmerkung beigefügt wurde. Im Übrigen lassen sich die historischen Angaben, die \textsc{Gauss} in den Anzeigen seiner arithmetischen Abhandlungen veröffentlicht hat, mithilfe des Nachlasses dahin ergänzen, dass die in den Artikeln $15$ bis $20$ der \emph{Theoria residuorum biquadrat.} enthaltenen Lehrsätze schon vor der Ausarbeitung der \emph{Theoria motus corporum coel.} niedergeschrieben wurden. Die in den Anzeigen erwähnten Untersuchungen über kubische Reste sind wohl nicht zur Ausarbeitung gelangt; aufgezeichnet sind davon die mit den Hilfsmitteln der Abhandlung \emph{Disquisitionum circa aequationes puras ulterior evolutio} geführten Beweise der Reziprozitätssätze für zwei Primzahlen, von denen eine reell ist.

Die Bruchstücke III bis VI bilden in der Handschrift besondere Hefte. Bei den ersten drei weist die Form der Schriftzüge auf eine Zeit hin, die nicht weit von derjenigen der Bruchstücke I und II entfernt liegt, während für das letzte, Nr. VI, ein erheblich späterer Zeitpunkt anzunehmen ist.

\par\medskip
[I.] Art. $10$. Die Anzahl der Gitterpunkte in $(z,z',z'+\zeta,z+\zeta)$ ergibt sich aus dem Satz: Sind $a$ und $b$ teilerfremd, so geht die von $\xi+\eta i$ nach $\xi+\eta i+a+bi$ gezogene Gerade durch einen Gitterpunkt, wenn der imaginäre Teil von $(\xi+\eta i)\pdot(-a+bi)$ eine ganze Zahl ist.

\SourceLeaf{386}

[I.] Art. $17$. Die erste Umformung des letzten $S$ im Ausdruck für $\Delta S$ erhält man, indem man das betreffende Flächenstück in drei Teile zerlegt, sodass jenes $S$ in
\[
\begin{aligned}
&S\bigl(\tfrac12m+\tfrac12-\tfrac12i,\ (\tfrac12+\tfrac12i)m,\\
&\qquad(\tfrac12+\tfrac12i)m+i,\ \tfrac12m+\tfrac12+\tfrac12i_{(+)}\bigr)\\
&-S\bigl(\tfrac12m,\ \tfrac12m+\tfrac12+\tfrac12i_{(+)},\ \tfrac12m+\tfrac12-\tfrac12i\bigr)\\
&-S\bigl((\tfrac12+\tfrac12i)m,\ \tfrac12m,\ \tfrac12m+\tfrac12-\tfrac12i\bigr)
\end{aligned}
\]
übergeht, dann die Gitterpunkte im ersten Flächenteil mithilfe des Satzes in Artikel $10$ auszählt und außerdem berücksichtigt, dass im zweiten Flächenteil kein Gitterpunkt liegt.

Die zweite Umformung erhält man, indem man die den Eckpunkten des dritten Flächenteils entsprechenden Größen mit $i$ multipliziert und um die ganze Zahl $(1-i)\frac{m-1}{2}$ vermehrt; schließlich erhält man die dritte, indem man die zuletzt entstandene Figur nach der Vorschrift von Artikel $16$ mit derjenigen vergleicht, die gegenüber ihr die Ortsverschiebung $\frac{-1}{1+i}$ hat.

[I.] Art. $[18.]$ Eine Erläuterung zum ersten Schema findet sich in den Artikeln $1$ bis $6$ des später niedergeschriebenen, hier mit [II.] bezeichneten Bruchstücks.

[I.] Art. $[18.]\ (2.)$ Die geometrische Deutung ergibt mithilfe der beiden Systeme von Gitterpunkten
\[
[-2iQ-\tfrac12,\tfrac12,iQ]=[-2iQ+\tfrac12,\tfrac12,iQ]=\mathrm{IV}^{*}
\]
und
\[
-[-2iQ+\tfrac12i,-\tfrac12i,iQ]=-[-2iQ-\tfrac12i,-\tfrac12i,iQ]=\mathrm{X}^{*}
\]
die Gleichungen
\[
\begin{aligned}
\mathrm{IV}-\mathrm{IV}^{*}+[-2iQ,\tfrac12i,-\tfrac12]
&=-\mathrm{X}+\mathrm{X}^{*}+[-iQ,\tfrac12i,-\tfrac12]\\
\mathrm{XIII}+\mathrm{IV}^{*}&=[-2iQ,1,iQ]=-I(-2iQ)+I(-iQ)\\
\mathrm{V}+\mathrm{X}^{*}&=-[-2iQ,-i,iQ]=R(-2iQ)-R(-iQ)
\end{aligned}
\]
wobei allgemein $Rx$ und $Ix$ die größten im reellen Teil und im Koeffizienten des imaginären Teils von $x$ enthaltenen ganzen Zahlen bezeichnen.

$[-2iQ,\tfrac12i,-\tfrac12]$ ist aber die Anzahl der Gitterpunkte in dem Quadrat, dessen Mittelpunkt bei $-2iQ-\tfrac14+\tfrac14i$ liegt und zwischen dessen Endpunkten die Ortsverschiebungen $\tfrac12$ und $\tfrac12i$ bestehen.

$[-iQ,\tfrac12i,-\tfrac12]$ oder $[-2iQ,-\tfrac12i,\tfrac12]$ ist die Anzahl der Gitterpunkte in einem gleichen Quadrat mit Mittelpunkt $-2iQ+\tfrac14-\tfrac14i$.

[I.] Art. $[18.]\ (3.)$ Mithilfe der Gitterpunkte
\[
[0,\tfrac12i,-iQ]=-[-\tfrac12i,-\tfrac12i,-iQ]=\mathrm{I}^{*}
\]
erhält man

\SourceLeaf{387}

\[
\begin{aligned}
\mathrm{VII}-\mathrm{XII}&=\mathrm{I}-\mathrm{I}^{*}+[-iQ,\tfrac12,\tfrac12i]-[0,\tfrac12,\tfrac12i]\\
&=\mathrm{I}-\mathrm{I}^{*}+[-2iQ,\tfrac12,\tfrac12i]\\
\mathrm{II}-\mathrm{I}^{*}&=[0,-i,-iQ]=-R0+R(-iQ)\\
\mathrm{IX}-\mathrm{I}&=[0,-1,-iQ]=I0-I(-iQ)\\
\mathrm{VI}-\mathrm{XII}&=-[\tfrac12i,-1,-iQ]=-I(\tfrac12i)+I(-iQ+\tfrac12i)\\
\mathrm{VIII}-\mathrm{XIV}&=-[\tfrac12,-1-i,\tfrac12im_{(-)}]=\frac{a-1}{2}+\frac b2
\end{aligned}
\]

[II.] Art. $10$. Es gilt
\[
\begin{aligned}
\Upsilon P'&\equiv-\Upsilon(-iP'+\tfrac12im),\\
\Upsilon P''&\equiv-1-\Upsilon(P''-\tfrac12im),\\
\Upsilon P'''&\equiv1+\Upsilon(-iP''')\ (\text{mod. }4)
\end{aligned}
\]
und $-iP'+\tfrac12im$, $P''-\tfrac12im$ sind die jeweils um $\tfrac12i$ vermehrten Gitterpunkte in I und VI.

[III.] Art. $6$. Die in der Anmerkung angegebenen Regeln zur Bestimmung von $\mathrm{Dec.}\frac Mm$ habe ich der vorliegenden Abhandlung aus einer anderen Stelle der Handschriften beigefügt. Die erste dieser beiden Regeln, die, wie leicht zu sehen ist, mit der zweiten übereinstimmt, folgt aus derjenigen von Artikel $6$, weil
\[
\begin{gathered}
k\equiv f\pdot i^{-n}\ (\text{mod. }m),\quad n=\Theta\frac{2fM}{m},\quad p=mm',\\
\left[\frac{2km'}p\right]^2\equiv\Theta\frac{2km'}p\ (\text{mod. }2+2i)
\end{gathered}
\]
gilt.

[III.] Art. $8$ enthält in der Handschrift ein Beispiel zu Artikel $7$, nämlich die Bestimmung des Dezidenten von $-1+2i$ für den Modul $-11+4i$.

[III.] Art. $10$. Zur Bemerkung ‘anders auszudrücken’ kann man Artikel $3$ des folgenden Bruchstücks [IV] vergleichen.

[IV.] Die Artikel $1$. $2$. $4$ enthalten in der Handschrift außer dem hier Abgedruckten noch die Anwendung auf die beiden Beispiele für $m=5+8i$, $M=9+4i$ und für $m=9+4i$, $M=5+8i$.
